% (c)~2014 Claudio Carboncini - claudio.carboncini@gmail.com
% (c)~2014 Dimitrios Vrettos - d.vrettos@gmail.com
% (c) 2015 Daniele Zambelli daniele.zambelli@gmail.com

\section{Esercizi}

\subsection{Esercizi dei singoli paragrafi}

\begin{comment}

\subsubsection*{\numnameref{sec:01_}}

\begin{esercizio}
\label{ese:D.19}
testo esercizio
\end{esercizio}

\begin{esercizio}\label{ese:03.1}
Consegna:
\begin{enumeratea}
\item  
\end{enumeratea}
\end{esercizio}

\subsection{Esercizi riepilogativi}

\begin{esercizio}
\label{ese:D.19}
testo esercizio
\end{esercizio}

\begin{esercizio}\label{ese:03.1}
Consegna:
\begin{enumeratea}
\item  
\end{enumeratea}
\end{esercizio}

\end{comment}

\subsection*{\numnameref{sec:sist2gr_sisteminumerci}}

\begin{esercizio}[*]
\label{ese:6.1}
Risolvere i seguenti sistemi di secondo grado.
\begin{htmulticols}{2}
\begin{enumeratea}
\item~\(\sistema{x^2+2y^2=3\\x+y=2}\)
\sol{\left(1;~1\right) \vee 
      \left(\dfrac 5 3;~\dfrac{1}{3}\right)}
\item~\(\sistema{4x^2+2y^2-6=0\\x=y}\)
\sol{\left(1;~1\right)\vee \left(-1;~-1\right)}
\item~\(\sistema{y^2-3y=2xy\\y=x-3}\)
\sol{\left(3;~0\right)\vee \left(-6;~-9\right)}
\item~\(\sistema{xy-x^2+2y^2=y-2x\\x+y=0}\)
\sol{\left(0;~0\right)}
\item~\(\sistema{{x+2y=-1}\\{x+5y^2=23}}\)
\sol{\left(-\dfrac{29} 5,\dfrac{12} 5\right);~   
            (3,-2)}
\item~\(\sistema{{x-5y=2}\\{x^2+2y^2=4}}\)
\sol{\left(-\dfrac{46}{27},-\dfrac{20}{27}\right);~
(2,0)}
\item~\(\sistema{3x-y=2\\x^2+2xy+y^2=0}\)
\sol{\left(\dfrac 1 2;~-\dfrac 1 2\right)}
\item~\(\sistema{x^2+y^2=1\\x+3y=10}\)
\sol{\emptyset}
\item~\(\sistema{x^2+y^2=2\\x+y=2}\)
\sol{\coppia{1}{1}}
\item~\(\sistema{3x+y=2\\x^2-y^2=1}\)
\sol{\emptyset}
%  \item~\(\sistema{x+y=0\\x^2+y^2-x-10=0}\)
% \sol{\left(-2;~2\right)\vee \left(\dfrac 5 2;~-\dfrac 5 
% 2\right)}
%  \item~\(\sistema{x^2-4y^2=0\\4x-7y=2}\)
% \sol{\left(4;~2\right)\vee \left(\dfrac 4{15};~-\dfrac 
% 2{15}\right)}
\end{enumeratea}
\end{htmulticols}
\end{esercizio}

\begin{esercizio}[*]
\label{ese:6.1}
Risolvere i seguenti sistemi di secondo grado.
\begin{enumeratea}
\item~\(\sistema{x^2-4xy+4y^2-1=0\\x=y+2}\)
\sol{\left(3;~1\right)\vee \left(5;~3\right)}
\item~\(\sistema{2x^2-6xy=x\\3x+5y=-2}\)
\sol{\left(0;~-\dfrac 2 5\right)\vee 
      \left(-\dfrac 1 4;~-\dfrac 1 4\right)}
\item~\(\sistema{x+y=1\\x^2+y^2-3x+2y=3}\)
\sol{\left(0;~1\right)\vee \left(\dfrac 7 2;~-\dfrac 5 
2\right)}
\item~\(\sistema{x^2+y^2=25\\4x-3y+7=0}\)
\sol{\left(-4;~-3\right) \vee 
      \left(\dfrac{44}{25};~\dfrac{117}{25}\right)}
\item~\(\sistema{x^2-4xy+4y^2-1=0\\x=y+2}\)
\sol{\left(3;~1\right)\vee \left(5;~-3\right)}
\item~\(\sistema{2x^2+xy-7x-2y=-6\\2x+y=3}\)
\sol{\left(1;~-3\right)\vee \left(-8;~-\dfrac{15} 2\right)}
%  \item~\(\sistema{x+y=0\\x^2+y^2-x-10=0}\)
% \sol{\left(-2;~2\right)\vee \left(\dfrac 5 2;~-\dfrac 5 
% 2\right)}
%  \item~\(\sistema{x^2-4y^2=0\\4x-7y=2}\)
% \sol{\left(4;~2\right)\vee \left(\dfrac 4{15};~-\dfrac 
% 2{15}\right)}
\end{enumeratea}
\end{esercizio}

\begin{esercizio}[*]
\label{ese:6.7}
Risolvere i seguenti sistemi di secondo grado.
\begin{enumeratea}
\item~\(\sistema{3x^2-4y^2-x=0\\x-2y=1}\)
\sol{\left(-1;~-1\right)\vee \left(\dfrac 1 2;~-\dfrac 1 
4\right)}

\item~\(\sistema{5x^2-y^2+4y-2x+2=0\\x-y=1}\)
\sol{\left(-\dfrac 3 2;~-\dfrac 5 2\right) \vee 
      \left(\dfrac 1 2;~-\dfrac 1 2\right)}

\item~\(\sistema{x+2y=3\\x^2-4xy+2y^2+x+y-1=0}\)
\sol{\left(1;~1\right)\vee \left(\dfrac{10} 7;
  \dfrac{11}{14}\right)}
\item~\(\sistema{x-2y-7=0\\x^2-xy=4}\)
\sol{\left(1;~-3\right)\vee \left(22;~\dfrac{15}{2}\right)}
% \sol{\forall (x,y)\in \R \times \R:\,y=-2x+3} 
% ????????
\item~\(\sistema{x^2+2y^2-3xy-x+2y-4=0\\2x-3y+4=0}\)
\sol{\left(4;~4\right)\vee \left(-5;~-2\right)}
\item~\(\sistema{x-2y=1\\x^2+y^2-2x=1}\)
\sol{\left(1+\dfrac{2\sqrt{10}} 5;~\dfrac{\sqrt{10}} 5\right)\vee 
\left(1-\dfrac{2\sqrt{10}} 5;~-\dfrac{\sqrt{10}} 5\right)}

\item~\(\sistema{x+y=1\\x^2+y^2-2xy-2y-2=0}\)
\sol{\left(\dfrac{1+\sqrt{13}} 4;~\dfrac{3-\sqrt{13}} 4\right)\vee 
\left(\dfrac{1-\sqrt{13}} 4;~\dfrac{3+\sqrt{13}} 4\right)}
%  
% \item~\(\sistema{9x^2-12xy+4y^2-2x+6y=8\\x-2y=2}\)
% \sol{\left(\dfrac{-9+\sqrt{241}} 
% 8;~\dfrac{-25+\sqrt{241}}{16}\right)\vee 
% \left(\dfrac{9-\sqrt{241}} 8;~\dfrac{-25-\sqrt{241}}{16}\right)}
%  \item~\(\sistema{3x+y=4\\x^2-y^2=1}\)
% \sol{\left(\dfrac{6-\sqrt 2} 4;~\dfrac{-2+3\sqrt 2} 4\right)\vee 
% \left(\dfrac{6+\sqrt 2} 4;~\dfrac{-2-3\sqrt 2} 4\right)}
%  \item~\(\sistema{\dfrac 1 2(2y-x)(y+x)-(x+y)^2+\dfrac 3 
% 2x(x+y+1)+2(y-1)=0\\\dfrac 2 3(x-3)^2+4\left(x-\dfrac 3 
% 2\right)=2({xy}+1)}\)
% \sol{\left(\dfrac{6-8\sqrt 
% 3}{13};~\dfrac{17+12\sqrt 3}{26}\right)\vee\left(\dfrac{6+ 8\sqrt 
% 3}{13};~\dfrac{17-12\sqrt 3}{26}\right)}
\end{enumeratea}
\end{esercizio}

\begin{esercizio}[*]
\label{ese:6.8}
Risolvere i seguenti sistemi, dopo aver eseguito la discussione sul 
parametro.
\begin{enumeratea}
\item~\(\sistema{x+y=3 \\x^2+y^2=k }\)
\sol{k\ge \dfrac 9 2.~\left(\dfrac{3-\sqrt{2k-9}} 2;~ 
\dfrac{3+\sqrt{2k-9}} 2\right)\vee \left(\dfrac{3+\sqrt{2k-9}} 2;~ 
\dfrac{3-\sqrt{2k-9}} 2\right)}
\item~\( \sistema{ky+2x=4\\xy=2}\)
\sol{...}
\item~\( \sistema{y=kx-1 \\y^2-kx^2+1=0}\)
\sol{...}
\item~\(\sistema{y=kx-2k \\x^2-2y-x=2}\)
\sol{\forall k\in \R:~(2;~ 0)\vee (2k-1;~ 2k^2-3k)}
\item~\(\sistema{y=x+k \\y=3x^2+2x}\)
\item~\( \sistema{y=-x+k \\x^2-y^2-1=0}\)
\item~\(\sistema{y+x-k=0 
\\{xy}+2{kx}-3{ky}-6k^2=0}\)
\item~\( \sistema{y-x+k=0 
\\y-x^2+4x-3=0}\)
\end{enumeratea}
\end{esercizio}

\begin{esercizio}[*]
\label{ese:6.10}
Trovare le soluzioni dei seguenti sistemi frazionari.
\begin{htmulticols}{2}
\begin{enumeratea}

\item~\(\sistema{x^2+y^2=4\\\dfrac{x+2y}{x-1}=2}\)



\item~\(\sistema{\dfrac{x+2y}{x-y}=4\\x^2+y^2+3x-2y=1}\)

\item~\(\sistema{\dfrac{2x+y}{x+2y}=3\\xy+3y=1}\)
\item~\(\sistema{\dfrac{3x-2y} 
x=\dfrac{1-x}{y-1}\\2x-y=1}\)
\end{enumeratea}
\end{htmulticols}
\end{esercizio}

\paragraph{6.9.} a)~\(k\ge -\dfrac 1{12}:~\left(\dfrac{-1-\sqrt{12k+1}} 
6;~ 
\dfrac{6k-1-\sqrt{12k+1}} 6\right) \vee \left(\dfrac{-1+\sqrt{12k+1}} 6;~ 
\dfrac{6k-1+\sqrt{12k+1}} 6\right)\),\protect\\
\quad d)~\(\forall k\in \R:~(3k;~ -2k)\)

\paragraph{6.10.} a)~\(x\neq 1:~\left(2;~0\right)\vee \left(-\dfrac 6 
5;~-\dfrac 8 
5\right)\),\quad b)~\(x\neq y:~\left(\dfrac 2 5;~\dfrac 1 5\right)\vee 
\left(-2;~-1\right)\),\quad c)~\(x\neq -2y:~\emptyset\),\protect\\
\quad d)~\(x\neq 0\wedge y\neq 1:~(4;~7)\)

\begin{esercizio}[*]
\label{ese:6.11}
Trovare le soluzioni dei seguenti sistemi frazionari.
\begin{htmulticols}{2}
\begin{enumeratea}
\item~\(\sistema{\dfrac{x+y}{x-2}=y+\dfrac 1 
3\\y=2x+2}\)

\item~\(\sistema{
\dfrac{2x+1}{y-2}=\dfrac{y-1}{x+1}\\2x+2y=3
}\)


\item~\(\sistema{\dfrac{y-1}{x+y}=x\\x-y=0}\)

\item~\(\sistema{\dfrac{x+1}{2y-1}=y\\2y-x=-4}\)
\end{enumeratea}
\end{htmulticols}
\end{esercizio}

\begin{esercizio}
\label{ese:6.12}
Risolvere i seguenti sistemi di secondo grado in tre incognite.
\begin{htmulticols}{2}
\begin{enumeratea}



\item~\(\sistema{x-3y-z=-4\\3x+2y+z=6\\4x^2+2xz+y^2=6}\)

\item~\(\sistema{x+y=5\\2x-y+3z=9\\x^2-y+z^2=1}\)

\item~\(\sistema{x-y+z=1\\2x-y+z=0\\x^2-y+z=3}\)


\item~\(\sistema{x-y+2z=3\\2x-2y+z=1\\x^2-y^2+z=12}\)
\end{enumeratea}
\end{htmulticols}
\end{esercizio}

\paragraph{6.11.} a)~\(x\neq 2:~\left(-1;~0\right)\vee \left(\dfrac{10} 
3;~\dfrac{26} 3\right)\),\quad b)~\(x\neq -1\wedge y\neq 2:~\left(-\dfrac 
5 
2;~4\right)\),\quad c)~\(x\neq -y:~\emptyset\),\protect\\
\quad d)~\(y\neq \dfrac 1 2:~\left(2;~-1\right)\vee \left(9;~\dfrac 5 
2\right)\)

\paragraph{6.12.} a)~\(\left(1;~2;~-1\right)\),\quad 
b)~\(\emptyset\),\quad 
c)~\(\forall z \in \R (-1;~z-2;~z)\),\quad d)~\(\left(-\dfrac{47} 
3;~-\dfrac{46} 
3;~\dfrac 5 3\right)\)

\begin{esercizio}[*]
\label{ese:6.13}
Risolvere i seguenti sistemi di secondo grado in tre incognite.
\begin{htmulticols}{2}
\begin{enumeratea}

\item~\(\sistema{2x-3y=-3\\5y+2z=1\\x^2+y^2+z^2=1}\)



\item~\(\sistema{x-2y+z=3\\x+2y+z=3\\x^2+y^2+z^2=29}\)

\item~\(\sistema{x+y-z=0\\x-y+3z=9\\x^2-y+z=12}\)


\item~\(\sistema{x-y=1\\x+y+z=0\\x^2+xy-z=0}\)


\item~\(\sistema{x-y-z=-1\\x+y+z=1\\x+y^2+z^2=32}\)
\end{enumeratea}
\end{htmulticols}
\end{esercizio}

\paragraph{6.13.} a)~\(\emptyset\),\quad b)~\(\left(5;~0;~-2\right)\vee 
\left(-2;~0;~5\right)\),\quad c)~\(\left(-4;~\dfrac{25} 2;~\dfrac{17} 
2\right)\vee 
\left(3;~-\dfrac 3 2;~\dfrac 3 2\right)\),\quad 
d)~\(\left(-1;~-2;~3\right)\vee 
\left(\dfrac 1 2;~-\dfrac 1 2;~0\right)\),\quad 
e)~\(\left(0;~\dfrac{3\sqrt 
7+1} 
2;~-\dfrac{3\sqrt 7-1} 2\right)\vee \left(0;~-\dfrac{3\sqrt 7-1} 
2;~\dfrac{3\sqrt 7+1} 
2\right)\)

\subsection*{5.2 - Equazioni riconducibili al prodotto di due o più 
fattori}

% \subsection*{6.2 - Sistemi simmetrici}
\subsection*{\numnameref{sec:eq2gr_sistemi_simmetrici}}

\begin{esercizio}[*]
\label{ese:6.14}
Risolvere i seguenti sistemi simmetrici di secondo grado.
\begin{htmulticols}{2}
\begin{enumeratea}
\item~\(\sistema{x+y=4\\{xy}=3}\)
\sol{(3;~1)\vee(1;~3)}
\item~\(\sistema{x+y=1\\{xy}=7 }\)
\sol{\emptyset}
\item~\(\sistema{x+y=5\\{xy}=6 }\)
\sol{(3;~2)\vee(2;~3)}
\item~\(\sistema{x+y=-5\\{xy}=-6 }\)
\sol{(1;~-6)\vee(-6;~1)}
\item~\(\sistema{x+y=3\\{xy}=2 }\)
\sol{(2;~1)\vee(1;~2)}
\item~\(\sistema{x+y=3\\{xy}=-4}\)
\sol{(4;~-1)\vee(-1;~4)}
\item~\(\sistema{x+y=-4\\{xy}=4 }\)
\sol{(-2;~-2)}
\item~\(\sistema{x+y=6\\{xy}=9 }\)
\sol{(3;~3)}
\item~\(\sistema{x+y=2\\{xy}=10 }\)
\sol{\emptyset}
\item~\(\sistema{x+y=7\\{xy}=12 }\)
\sol{(4;~3)\vee(3;~4)}
\item~\(\sistema{x+y=-1\\{xy}=2 }\)
\sol{\emptyset}
\item~\(\sistema{x+y=12\\{xy}=-13 }\)
\sol{(13;~-1)\vee(-1;~13)}
%  \item~\(\sistema{
%   x+y=\dfrac 6 5\\{xy}=\dfrac 9{25}
%  }\)
%   \sol{\left(\dfrac 3 5;~\dfrac 3 5\right)}
%  \item~\(\sistema{x+y=4\\{xy}=50 }\)
%   \sol{\emptyset}
%  \item~\(\sistema{x+y=-5\\{xy}=-14 }\)
%   \sol{(2;~-7)\vee(-7;~2)}
%  \item~\(\sistema{x+y=5\\{xy}=-14}\)
%   \sol{(7;~-2)\vee(-2;~7)}
%  \item~\(\sistema{x+y=\dfrac 1 4\\{xy}=-\dfrac 3 
% 8}\)
%   \sol{\left(\dfrac 3 4;~-\dfrac 1 2\right)\vee
%          \left(-\dfrac 1 2;~\dfrac 3 4\right)}
%  \item~\(\sistema{x+y=2\\{xy}=-10}\)
%   \sol{\left(1+\sqrt{11};~1-\sqrt{11}\right)\vee
%          \left(1-\sqrt{11};~1+\sqrt{11}\right)}
\end{enumeratea}
\end{htmulticols}
\end{esercizio}

\begin{esercizio}[*]
\label{ese:6.18}
Risolvere i seguenti sistemi simmetrici di secondo grado.
\begin{htmulticols}{2}
\begin{enumeratea}
\item~\(\sistema{x+y=4\\{xy}=0 }\)
\sol{}
\item~\(\sistema{x+y=\dfrac 5 2\\{xy}=-\dfrac 7 
2}\)
\item~\(\sistema{x+y=-5\\{xy}=2 }\)
\item~\(\sistema{x+y=\dfrac 4 3\\{xy}=-\dfrac 1 2 
}\)
\end{enumeratea}
\end{htmulticols}
\end{esercizio}

\begin{esercizio}[*]
\label{ese:6.19}
Risolvere i seguenti sistemi simmetrici di secondo grado.
\begin{htmulticols}{2}
\begin{enumeratea}
\item~\(\sistema{x+y=\dfrac 5 2\\{xy}=-\dfrac 9 
2}\)
\item~\(\sistema{x+y=2\\{xy}=-\dfrac 1 
3}\)
\item~\(\sistema{x+y=1\\{xy}=-3}\)
\item~\(\sistema{x+y=4\\{xy}=-50}\)
\end{enumeratea}
\end{htmulticols}
\end{esercizio}

\paragraph{6.18.} a)~\((0;~4)\vee(4;~0)\),\;~ b)~\(\left(\dfrac 7 
2;~-1\right)\vee\left(-1;~\dfrac 7 2\right)\),\;~ 
c)~\(\left(\dfrac{-5+\sqrt{17}} 
2;~\dfrac{-5-\sqrt{17}} 2\right)\vee \left(\dfrac{-5-\sqrt{17}} 
2\\y=\dfrac{-5+\sqrt{17}} 2\right)\),\quad d)~\(\left(\dfrac{4+\sqrt{34}} 
6;~\dfrac{4-\sqrt{34}} 6\right)\vee \left(\dfrac{4-\sqrt{34}} 
6\\y=\dfrac{4+\sqrt{34}} 6\right)\)

\paragraph{6.19.} a)~\(\left(\dfrac{5+\sqrt{97}} 4\\y=\dfrac{5-\sqrt{97}} 
4\right)\vee \left(\dfrac{5-\sqrt{97}} 4\\y=\dfrac{5+\sqrt{97}} 
4\right)\),\quad 
b)~\(\left(\dfrac{3+{2\sqrt 3}} 3;~\dfrac {3-{2\sqrt 3}} 3\right)\vee 
\left(\dfrac 
{3-{2\sqrt 3}} 3;~\dfrac{3+{2\sqrt 3}} 3\right)\),\quad 
c)~\(\left(\dfrac{1+\sqrt{13}} 2;~\dfrac{1-\sqrt{13}} 2\right)\vee 
\left(\dfrac{1-\sqrt{13}} 2;~\dfrac{1+\sqrt{13}} 2\right)\),\quad 
d)~\(\left(2+3\sqrt 
6;~2-3\sqrt 6\right)\vee \left(2-3\sqrt 6;~2+3\sqrt 6\right)\)

\begin{esercizio}[*]
\label{ese:6.20}
Risolvere i seguenti sistemi riconducibili al sistema simmetrico 
fondamentale.
\begin{htmulticols}{2}
\begin{enumeratea}
\item~\(\sistema{x+y=1 \\x^2+y^2=1}\)
\item~\(\sistema{x+y=2 \\x^2+y^2=2}\)
\item~\(\sistema{x+y=3 \\x^2+y^2=5}\)
\item~\(\sistema{x+y=2 \\x^2+y^2+x+y=1}\)
\item~\(\sistema{x+y=4\\x^2+y^2=8}\)
\item~\(\sistema{x+y=2 \\x^2+y^2-3xy=4}\)
\end{enumeratea}
\end{htmulticols}
\end{esercizio}

\begin{esercizio}[*]
\label{ese:6.21}
Risolvere i seguenti sistemi riconducibili al sistema simmetrico 
fondamentale.
\begin{htmulticols}{2}
\begin{enumeratea}

\item~\(\sistema{{x+y=-12}\\{x^2+y^2=72}}\)

\item~\(\sistema{{2x+2y=-2}\\{(y-x)^2-{xy}=101}}\)



\item~\(\sistema{{-4x-4y=-44}\\{2x^2+2y^2-3{xy}=74}}\)
\item~\(\sistema{x+y=3 
\\x^2+y^2-4x-4y=5}\)
\end{enumeratea}
\end{htmulticols}
\end{esercizio}

\begin{esercizio}[*]
\label{ese:6.22}
Risolvere i seguenti sistemi riconducibili al sistema simmetrico 
fondamentale.
\begin{htmulticols}{2}
\begin{enumeratea}
\item~\(\sistema{x+y=7 \\x^2+y^2=29}\)
\item~\(\sistema{2x+2y=-2\\4x^2+4y^2=52}\)
\item~\(\sistema{\dfrac{x+y} 2=\dfrac 3 
4\\3x^2+3y^2=\dfrac{15} 
4}\)
\item~\(\sistema{x+y=-3 
\\x^2+y^2-5xy=37}\)
\end{enumeratea}
\end{htmulticols}
\end{esercizio}

\begin{esercizio}[*]
\label{ese:6.23}
Risolvere i seguenti sistemi riconducibili al sistema simmetrico 
fondamentale.
\begin{htmulticols}{2}
\begin{enumeratea}
\item~\(\sistema{x+y=-6 \\x^2+y^2-xy=84}\)
\item~\(\sistema{x+y=-5 
\\x^2+y^2-4xy+5x+5y=36}\)
\item~\(\sistema{x+y=-7 
\\x^2+y^2-6xy-3x-3y=44}\)
\item~\(\sistema{x^2+y^2=-1\\x+y=6 }\)
\end{enumeratea}
\end{htmulticols}
\end{esercizio}

\begin{esercizio}[*]
\label{ese:6.24}
Risolvere i seguenti sistemi riconducibili al sistema simmetrico 
fondamentale.
\begin{htmulticols}{2}
\begin{enumeratea}
\item~\(\sistema{x^2+y^2=1\\x+y=-7}\)
\item~\(\sistema{x^2+y^2=18\\x+y=6 }\)
\item~\(\sistema{x^2+y^2-4xy-6x-6y=1\\x+y=1 
}\)
\item~\(\sistema{x^2+y^2=8\\x+y=3}\)
\end{enumeratea}
\end{htmulticols}
\end{esercizio}

\begin{esercizio}[*]
\label{ese:6.25}
Risolvere i seguenti sistemi riconducibili a sistemi simmetrici.
\begin{htmulticols}{3}
\begin{enumeratea}
\item~\(\sistema{{x-y=1}\\{x^2+y^2=5}}\)
\item~\(\sistema{{\dfrac 1 x+\dfrac 1 y=-12}\\{{xy}=\dfrac 
1{35}}}\)
\item~\(\sistema{{-2x+y=3}\\{{xy}=1}}\)
\end{enumeratea}
\end{htmulticols}
\end{esercizio}

\begin{esercizio}[*]
\label{ese:6.26}
Risolvere i seguenti sistemi di grado superiore al secondo.
\begin{htmulticols}{2}
\begin{enumeratea}
\item~\(\sistema{{x+y=-1}\\{x^3+y^3=-1}}\)

\item~\(\sistema{{{xy}=-2}\\{x^2+y^2=13}}\)


\item~\(\sistema{{x+y=-2}\\{x^3+y^3-{xy}=-5}}\)
\item~\(\sistema{{x+y=8}\\{x^3+y^3=152}}\)
\end{enumeratea}
\end{htmulticols}
\end{esercizio}

\begin{esercizio}[*]
\label{ese:6.27}
Risolvere i seguenti sistemi di grado superiore al secondo.
\begin{htmulticols}{2}
\begin{enumeratea}
\item~\(\sistema{x^3+y^3=9\\x+y=3}\)
\item~\(\sistema{x^3+y^3=-342\\x+y=-6}\)

\item~\(\sistema{{x^3-y^3=351}{{xy}=-14}}\)
\item~\(\sistema{x^3+y^3=35\\x+y=5}\)
\end{enumeratea}
\end{htmulticols}
\end{esercizio}

\begin{esercizio}[*]
\label{ese:6.28}
Risolvere i seguenti sistemi di grado superiore al secondo.
\begin{htmulticols}{2}
\begin{enumeratea}
\item~\(\sistema{x^4+y^4=2\\x+y=0}\)
\item~\(\sistema{x^4+y^4=17\\x+y=-3}\)
\item~\(\sistema{x^3+y^3=-35\\xy=6}\)
\item~\(\sistema{x^3+y^3=-26\\xy=-3}\)
\end{enumeratea}
\end{htmulticols}
\end{esercizio}

\begin{esercizio}[*]
\label{ese:6.29}
Risolvere i seguenti sistemi di grado superiore al secondo.
\begin{htmulticols}{2}
\begin{enumeratea}
\item~\(\sistema{{x+y=3}\\{x^4+y^4=17}}\)

\item~\(\sistema{{x+y=-1}\\{8x^4+8y^4=41}}\)
\item~\(\sistema{{x+y=3}\\{x^4+y^4=2}}\)
\item~\(\sistema{{x+y=5}\\{x^4+y^4=257}}\)
\end{enumeratea}
\end{htmulticols}
\end{esercizio}

\begin{esercizio}[*]
\label{ese:6.30}
Risolvere i seguenti sistemi di grado superiore al secondo.
\begin{htmulticols}{2}
\begin{enumeratea}
\item~\(\sistema{x^4+y^4=2\\xy=1}\)
\item~\(\sistema{x^4+y^4=17\\xy=-2}\)

\item~\(\sistema{{x+y=-1}\\{x^5+y^5=-211}}\)
\item~\(\sistema{x^5+y^5=64\\x+y=4}\)
\end{enumeratea}
\end{htmulticols}
\end{esercizio}

\begin{esercizio}[*]
\label{ese:6.31}
Risolvere i seguenti sistemi di grado superiore al secondo.
\begin{htmulticols}{2}
\begin{enumeratea}
\item~\(\sistema{x^5+y^5=-2882\\x+y=-2}\)
\item~\(\sistema{x^5+y^5=2\\x+y=0}\)
\item~\(\sistema{x^5+y^5=31\\xy=-2}\)
\item~\(\sistema{x^4+y^4=337\\xy=12}\)
\end{enumeratea}
\end{htmulticols}
\end{esercizio}

\begin{esercizio}[*]
\label{ese:6.32}
Risolvere i seguenti sistemi di grado superiore al secondo.
\begin{htmulticols}{2}
\begin{enumeratea}
\item~\(\sistema{x^3+y^3=\dfrac{511} 
8\\xy=-2}\)
\item~\(\sistema{x^2+y^2=5\\xy=2 }\)
\item~\(\sistema{x^2+y^2=34\\xy=15 }\)
\item~\(\sistema{xy=1 \\x^2+y^2+3xy=5}\)
\end{enumeratea}
\end{htmulticols}
\end{esercizio}

\begin{esercizio}[*]
\label{ese:6.33}
Risolvere i seguenti sistemi di grado superiore al secondo.
\begin{htmulticols}{2}
\begin{enumeratea}
\item~\(\sistema{xy=12 \\x^2+y^2=25}\)
\item~\(\sistema{xy=1 \\x^2+y^2-4xy=-2}\)
\item~\(\sistema{x^2+y^2=5\\xy=3 }\)
\item~\(\sistema{x^2+y^2=18\\xy=9 }\)
\end{enumeratea}
\end{htmulticols}
\end{esercizio}

\begin{esercizio}[*]
\label{ese:6.34}
Risolvere i seguenti sistemi di grado superiore al secondo.
\begin{htmulticols}{2}
\begin{enumeratea}
\item~\(\sistema{x^2+y^2+3xy=10\\xy=6 }\)
\item~\(\sistema{x^2+y^2+5xy-2x-2y=3\\xy=1 
}\)
\item~\(\sistema{x^2+y^2-6xy+3x+3y=2\\xy=2 
}\)
\item~\(\sistema{x^2+y^2=8\\xy=-3}\)
\end{enumeratea}
\end{htmulticols}
\end{esercizio}

\begin{esercizio}[*]
\label{ese:6.35}
Risolvere i seguenti sistemi di grado superiore al secondo.
\begin{enumeratea}
\item~\(\sistema{x^2+y^2+5xy+x+y=-6\\xy=-2 
}\)
\item~\(\sistema{{x+y=-\dfrac 1 
3}\\{x^5+y^5=-\dfrac{31}{243}}}\)
\item~\(\sistema{x^2+y^2+5xy+x+y=-\dfrac{25} 4\\xy=-2 
}\)
\item~\(\sistema{{x+y=1}\\{x^5+y^5=-2}}\)


\item~\(\sistema{{x+y=1}\\{x^5+y^5+7{xy}=17}}\)
\end{enumeratea}
\end{esercizio}

% \subsection*{6.3 - Sistemi omogenei di quarto grado}
% 
% \begin{esercizio}[*]
% \label{ese:6.36}
% Risolvi i seguenti sistemi omogenei.
% \begin{htmulticols}{2}
% \begin{enumeratea}
% 
% \item~\(\sistema{x^2-2xy+y^2=0\\x^2+3xy-2y^2=0}\)
% 
% \item~\(\sistema{3x^2-2xy-y^2=0\\2x^2+xy-3y^2=0}\)
% \item~\(\sistema{x^2-6{xy}+8y^2=0 \\x^2+4{xy}-5y^2=0 
% }\)
% 
% \item~\(\sistema{2x^2+xy-y^2=0\\4x^2-2xy-6y^2=0}\)
% \end{enumeratea}
% \end{htmulticols}
% \end{esercizio}
% 
% \paragraph{6.20.} a)~\((1,0)\vee(0,1)\),\quad b)~\((1,1)\),\quad 
% c)~\((1,2)\vee(2,1)\),\quad d)~\(\emptyset\),\quad e)~\((2,2)\),\quad 
% f)~\((0,2)\vee(2,0)\)
% 
% \paragraph{6.21.} a)~\((-6,-6)\),\quad b)~\((-5,4)\vee(4,-5)\),\quad 
% c)~\((3,8)\vee(8,3)\),\quad d)~\((-1,4)\vee(4,-1)\)
% 
% \paragraph{6.22.} a)~\((2,5)\vee(5,2)\),\quad 
% b)~\((-3,2)\vee(2,-3)\),\quad 
% c)~\((\dfrac 1 2;~1)\vee(1;~\dfrac 1 2)\),\quad d)~\((-4;~1)\vee(1;~-4)\)
% 
% \paragraph{6.23.} a)~\((-8;~2)\vee(2;~-8)\),\quad 
% b)~\((-6;~1)\vee(1;~-6)\),\quad 
% c)~\(\left(-\dfrac 1 2;~-\dfrac{13} 2\right)\vee \left(-\dfrac{13} 
% 2;~-\dfrac 1 
% 2\right)\),\quad d)~\(\emptyset\)
% 
% \paragraph{6.24.} a)~\(\emptyset\),\;~ b)~\((3;~3)\),\;~ 
% c)~\(\left(\dfrac{1+\sqrt 5} 
% 2;~\dfrac{1-\sqrt 5} 2\right)\vee \left(\dfrac{1-\sqrt 5} 
% 2;~\dfrac{1+\sqrt 5} 
% 2\right)\),\;~ d)~\(\left(\dfrac{3+\sqrt 7} 2;~\dfrac{3-\sqrt 7} 
% 2\right)\vee 
% \left(\dfrac{3-\sqrt 7} 2;~\dfrac{3+\sqrt 7} 2\right)\)
% 
% \paragraph{6.25.} a)~\( (-1;~ -2)\vee(2;~1)\),~b)~\(\left(-\dfrac 1 
% 7;~-\dfrac 1 
% 5\right)\vee\left(-\dfrac 1 5;~-\dfrac 1 
% 7\right)\),~c)~\(\left(\dfrac{-3-\sqrt{17}} 
% 4;~{\dfrac{3-\sqrt{17}} 2}\right)\vee \left(\dfrac{-3+\sqrt{17}} 
% 4;~{\dfrac{3+\sqrt{17}} 2}\right)\)
% 
% \paragraph{6.26.} a)~\((-1;~0)\vee(0;~-1)\),\quad 
% b)~\(\left(\dfrac{-3-\sqrt{17}} 
% 2;~\dfrac{-3+\sqrt{17}} 2\right)\vee\left(\dfrac{-3+\sqrt{17}} 
% 2;~\dfrac{-3-\sqrt{17}} 2\right)\),\protect\\
% \quad c)~\(\left(\dfrac{-5-\sqrt{10}} 5;~\dfrac{-5+\sqrt{10}} 
% 5\right)\vee 
% \left(\dfrac{-5+\sqrt{10}} 5;~\dfrac{-5-\sqrt{10}} 5\right)\),\quad 
% d)~\((3;~5)\vee(5;~3)\)
% 
% \paragraph{6.27.} a)~\((1;~2)\vee(2;~1)\),\quad 
% b)~\((-7;~1)\vee(1;~-7)\),\quad 
% c)~\((2;~-7)\vee(7;~-2)\),\quad d)~\((2;~3)\vee(3;~2)\)
% 
% \paragraph{6.28.} 
% a)~\((-1;~1)\vee(1;~-1)\),\;~b)~\((-2;~-1)\vee(-1;~-2)\),\;~ 
% c)~\((-3;~-2)\vee(-2;~-3)\),\;~ d)~\((-3;~1)\vee(1;~-3)\)
% 
% \paragraph{6.29.} a)~\((1;~2)\vee(2;~1)\),\quad b)~\(\left(-\dfrac 3 
% 2;~\dfrac 1 
% 2\right)\vee\left(\dfrac 1 2;~-\dfrac 3 2\right)\),\quad 
% c)~\(\emptyset\),\quad 
% d)~\((1;~4)\vee(4;~1)\)
% 
% \paragraph{6.30.} a)~\((1;~1)\),\quad 
% b)~\((1;~-2)\vee(-2;~1)\vee(-1;~2)\vee(2;~-1)\),\quad 
% c)~\((-3;~2)\vee(2;~-3)\),\quad 
% d)~\((2;~2)\)
% 
% \paragraph{6.31.} a)~\((-5;~3)\vee(3;~-5)\),\quad b)~\(\emptyset\),\quad 
% c)~\((-1;~2)\vee(2;~-1)\),\quad 
% d)~\((-4;~-3)\vee(-3;~-4)\vee(3;~4)\vee(4;~3)\)
% 
% \paragraph{6.32.} a)~\(\left(-\dfrac 1 2;~4\right)\vee\left(4;~-\dfrac 1 
% 2\right)\),\quad b)~\((-2;~-1)\vee(-1;~-2)\vee(1;~2)\vee(2;~1)\),\quad 
% c)~\((-5;~-3)\vee(-3;~-5)\vee(3;~5)\vee(5;~3)\),\quad 
% d)~\((-1;~-1)\vee(1;~1)\)
% 
% \paragraph{6.33.} a)~\((-4;~-3)\vee(-3;~-4)\vee(3;~4)\vee(4;~3)\),\quad 
% b)~\((-1;~-1)\vee(1;~1)\),\quad c)~\(\emptyset\),\quad 
% d)~\((-3;~-3)\vee(3;~3)\)
% 
% \paragraph{6.34.} a)~\(\emptyset\),\quad b)~\((1;~1)\),\quad 
% c)~\((-3-\sqrt 
% 7;~-3+\sqrt 
% 7)\vee(-3+\sqrt 7;~-3-\sqrt 7)\vee(1;~2)\vee(2;~1)\),\protect\\
% \quad d)~\(\left(\dfrac{\sqrt{14}+\sqrt 2} 2;~\dfrac{\sqrt 2-\sqrt{14}} 
% 2\right)\vee\left(\dfrac{\sqrt{14}-\sqrt 2} 2;~\dfrac{\sqrt 2+\sqrt{14}} 
% 2\right)\vee\left(\dfrac{\sqrt{14}-\sqrt 2} 2;~-\dfrac{\sqrt 2+\sqrt{14}} 
% 2\right)\vee\left(\dfrac{\sqrt{14}+\sqrt 2} 2;~-\dfrac{\sqrt 2-\sqrt{14}} 
% 2\right)\)
% 
% \paragraph{6.35.} a)~\((-2;~1)\vee(1;~-2)\vee(-\sqrt 
% {2};~\sqrt{2})\vee(\sqrt 
% {2};~-\sqrt{2})\),\quad b)~\(\left(-\dfrac 2 3;~\dfrac 1 
% 3\right)\vee\left(\dfrac 1 
% 3;~-\dfrac 2 3\right)\),\protect\\
% \quad c)~\(\left(\dfrac{-1+\sqrt{33}} 4;~\dfrac{-1-\sqrt{33}} 
% 4\right)\vee 
% \left(\dfrac{-1-\sqrt{33}} 4;~\dfrac{-1+\sqrt{33}} 4\right)\),\quad 
% d)~\(\emptyset\),\quad e)~\((-1;~2)\vee(2;~-1)\)
% 
% \paragraph{6.36.} a)~\((0;~0)\),\quad b)~\((t;~t)\),\quad 
% c)~\((0;~0)\),\quad 
% d)~\((t;~-t)\)
% 
% \begin{esercizio}[*]
% \label{ese:6.37}
% Risolvi i seguenti sistemi omogenei.
% \begin{htmulticols}{2}
% \begin{enumeratea}
% 
% \item~\(\sistema{x^2-5xy+6y^2=0\\x^2-4xy+4y^2=0}\)
% %  
% \item~\(\sistema{x^2-5xy+6y^2=0\\x^2+2xy-8y^2=0}\)
% %  
% \item~\(\sistema{x^2+xy-2y^2=0\\x^2+5xy+6y^2=0}\)
% %  
% % 
% \item~\(\sistema{x^2+7xy+12y^2=0\\2x^2+xy+6y^2=0}\)
% \end{enumeratea}
% \end{htmulticols}
% \end{esercizio}
% 
% \begin{esercizio}[*]
% \label{ese:6.38}
% Risolvi i seguenti sistemi omogenei.
% \begin{htmulticols}{2}
% \begin{enumeratea}
% 
% 
% \item~\(\sistema{x^2+6xy+8y^2=0\\2x^2+12xy+16y^2=0}\)
% \item~\(\sistema{-4x^2-7{xy}+2y^2=0 \\12x^2+21{xy}-6y^2=0 
% }\)
% 
% \item~\(\sistema{x^2+2xy+y^2=0\\x^2+3xy+2y^2=0}\)
% 
% 
% \item~\(\sistema{x^2+4xy=0\\x^2+2xy-4y^2-4=0}\)
% \end{enumeratea}
% \end{htmulticols}
% \end{esercizio}
% 
% \begin{esercizio}[*]
% \label{ese:6.39}
% Risolvi i seguenti sistemi omogenei.
% \begin{htmulticols}{2}
% \begin{enumeratea}
% 
% \item~\(\sistema{x^2-8xy+15y^2=0\\x^2-2xy+y^2=1}\)
% \item~\(\sistema{4x^2-y^2=0\\x^2-y^2=-3}\)
% 
% \item~\(\sistema{x^2+3xy+2y^2=0\\x^2-3xy-y^2=3}\)
% 
% 
% \item~\(\sistema{x^2-4xy+4y^2=0\\2x^2-y^2=-1}\)
% \end{enumeratea}
% \end{htmulticols}
% \end{esercizio}
% 
% \begin{esercizio}[*]
% \label{ese:6.40}
% Risolvi i seguenti sistemi omogenei.
% \begin{htmulticols}{2}
% \begin{enumeratea}
% 
% \item~\(\sistema{6x^2+5xy+y^2=12\\x^2+4xy+y^2=6}\)
% 
% \item~\(\sistema{x^2-xy-2y^2=0\\x^2-4xy+y^2=6}\)
% 
% \item~\(\sistema{x^2+y^2=3\\x^2-xy+y^2=3}\)
% 
% \item~\(\sistema{x^2-3xy+5y^2=1\\x^2+xy+y^2=1}\)
% \end{enumeratea}
% \end{htmulticols}
% \end{esercizio}
% 
% \begin{esercizio}[*]
% \label{ese:6.41}
% Risolvi i seguenti sistemi omogenei.
% \begin{htmulticols}{2}
% \begin{enumeratea}
% 
% 
% \item~\(\sistema{x^2+y^2=5\\x^2-3xy+y^2=11}\)
% 
% 
% \item~\(\sistema{x^2+5xy+4y^2=10\\x^2-2xy-3y^2=-11}\)
% \item~\(\sistema{4x^2-xy-y^2=-\dfrac 1 
% 2\\x^2+2xy-y^2=\dfrac 
% 1 
% 4}\)
% 
% \item~\(\sistema{x^2-xy-8y^2=-8\\x^2-2y^2-xy=16}\)
% \end{enumeratea}
% \end{htmulticols}
% \end{esercizio}
% 
% \begin{esercizio}
% \label{ese:6.42}
% Risolvi i seguenti sistemi omogenei.
% \begin{htmulticols}{2}
% \begin{enumeratea}
% 
% 
% \item~\(\sistema{x^2-6xy-y^2=10\\x^2+xy=-2}\)
% 
% \item~\(\sistema{4x^2-3xy+y^2=32\\x^2+3y^2-9xy=85}\)
% 
% \item~\(\sistema{x^2+3xy+2y^2=8\\3x^2-y^2+xy=-4}\)
% 
% 
% 
% \item~\(\sistema{x^2+5xy-7y^2=-121\\3xy-3x^2-y^2=-7}\)
% \end{enumeratea}
% \end{htmulticols}
% \end{esercizio}
% 
% \paragraph{6.37.} a)~\((2t;~t)\),\quad b)~\((2t;~t)\),\quad 
% c)~\((-2t;~t)\),\quad 
% d)~\((0;~0)\)
% 
% \paragraph{6.38.} a)~\((-4t;~t)\vee(-2t;~t)\),\quad 
% b)~\((k;~4k)\vee(k;~-\dfrac{1 
% 2}k)\),\quad c)~\((-t;~t)\),\quad d)~\((-4;~1)\vee(4;~-1)\)
% 
% \paragraph{6.39.} a)~\(\left(-\dfrac 3 2;~-\dfrac 1 
% 2\right)\vee\left(\dfrac 3 
% 2;~\dfrac 1 2\right)\vee\left(-\dfrac 5 4;~-\dfrac 1 
% 4\right)\vee\left(\dfrac 5 
% 4;~\dfrac 1 4\right)\),\quad b)~\((1;~2)\vee(-1;~-2)\vee (-1;~2)\vee 
% (1;~-2)\),\protect\\
% \quad c)~\((-1;~1)\vee (1;~-1)\left(-\dfrac{2\sqrt 3} 3;~\dfrac{\sqrt 3} 
% 3\right)\vee\left(\dfrac{2\sqrt 3} 3;~-\dfrac{\sqrt 3} 3\right)\),\quad 
% d)~\(\emptyset\)
% 
% \paragraph{6.40.} a)~\((1;~1)\vee(-1;~-1)\vee\left(\sqrt 6;~-4\sqrt 
% 6\right)\vee\left(-\sqrt 6;~4\sqrt 6\right)\),\quad 
% b)~\((1;~-1)\vee(-1;~1)\),\protect\\
% \quad c)~\((\sqrt 3;~0)\vee (-\sqrt 3;~0)\vee (0;~\sqrt 3)\vee (0;~-\sqrt 
% 3)\),\quad 
% d)~\((1;~0)\vee (-1;~0)\vee\left(\dfrac{\sqrt 3} 3;~\dfrac{\sqrt 3} 
% 3\right)\vee\left(-\dfrac{\sqrt 3} 3;~-\dfrac{\sqrt 3} 3\right)\)
% 
% \paragraph{6.41.} a)~\((1;~-2)\vee(-1;~2)\vee (-2;~1)\vee (2;~-1)\),\quad 
% b)~\((2;~-3)\vee(-2;~3)\),\protect\\
% \quad c)~\(\left(\dfrac 1 2;~1\right)\vee \left(-\dfrac 1 
% 2;~-1\right)\),\quad 
% d)~\((4;~-2)\vee(-4;~2)\vee (6;~2)\vee (-6;~-2)\)
% 
% \paragraph{6.42.} a)~\((-1;~3)\vee(1;~-3)\),\quad 
% b)~\((1;~-4)\vee(-1;~4)\vee 
% (-1;~-7)\vee (1;~7)\),\protect\\
% \quad c)~\((0;~2)\vee (0;~-2)\vee \left(\dfrac{10} 3;~-\dfrac{14} 
% 3\right)\vee 
% \left(-\dfrac{10} 3;~\dfrac{14} 3\right)\),\quad d)~\((2;~5)\vee 
% (-2;~-5)\vee 
% \left(-\dfrac{18} 7;~-\dfrac{37} 7\right)\vee \left(\dfrac{18} 
% 7;~\dfrac{37} 7\right)\)

\begin{esercizio}[*]
\label{ese:6.43}
Risolvi i seguenti sistemi particolari.
\begin{htmulticols}{2}
\begin{enumeratea}



\item~\(\sistema{x^2-5xy-3y^2=27\\-2x^2-2y^2+4xy=-50}\)

\item~\(\sistema{9x^2+5y^2=-3\\x^2+4xy-3y^2=8}\)



\item~\(\sistema{2x^2-4xy-3y^2=18\\xy-2x^2+3y^2=-18}\)
\item~\(\sistema{x^2+2xy=-\dfrac 7 
4\\x^2-4xy+4y^2=\dfrac{81} 
4}\)
\end{enumeratea}
\end{htmulticols}
\end{esercizio}

\begin{esercizio}[*]
\label{ese:6.44}
Risolvi i seguenti sistemi particolari.
\begin{htmulticols}{2}
\begin{enumeratea}



\item~\(\sistema{x^2+4xy+4y^2-16=0\\x^2-xy+4y^2-6=0}\)

\item~\(\sistema{x^2-2xy+y^2-1=0\\x^2-2xy-y^2=1}\)
\end{enumeratea}
\end{htmulticols}
\end{esercizio}

\begin{esercizio}
\label{ese:6.45}
Risolvi i seguenti sistemi particolari.
\begin{htmulticols}{2}
\begin{enumeratea}
\item~\(\sistema{x^2-y^2=0\\2x+y=3}\)

\item~\(\sistema{(x-2y)(x+y-2)=0\\3x+6y=3}\)

\item~\(\sistema{(x+y-1)(x-y+1)=0\\x-2y=1}\)

\item~\(\sistema{(x-3y)(x+5y-2)=0\\)x-2)(x-y+4)=0}\)
\end{enumeratea}
\end{htmulticols}
\end{esercizio}

\begin{esercizio}[*]
\label{ese:6.46}
Risolvi i seguenti sistemi particolari.
\begin{htmulticols}{2}
\begin{enumeratea}

\item~\(\sistema{(x^2-3x+2)(x+y)=0\\x-y=2}\)

\item~\(\sistema{(x-y)(x+y+1)(2x-y-1)=0\\)x-3y-3)(x+y-2)=0
}\)
\item~\(\sistema{(4x^2-9y^2)(x^2-2xy+y^2-9)=0 \\2x-y=2 
}\)


\item~\(\sistema{x^2+6xy+9y^2-4=0\\)x^2-y^2)(2x-y-4)=0}\)
\end{enumeratea}
\end{htmulticols}
\end{esercizio}

\begin{esercizio}[*]
\label{ese:6.47}
Risolvi i seguenti sistemi particolari.
\begin{enumeratea}

\item~\(\sistema{x^2-2xy-8y^2=0\\)x+y)(x-3)=0}\)
\item~\(\sistema{(2x^2-3xy+y^2)(x-y-1)=0 
\\)x^2-4xy+3y^2)(12x^2-xy-y^2)=0 }\)
\item~\(\sistema{(x-2y-2)(x^2-9y^2)=0 
\\)4x^2-4xy+y^2)(y+2)(x-y)=0 
}\)
\item~\(\sistema{x^4-y^4=0 \\x^2-(y^2-6y+9)=0 
}\)
\end{enumeratea}
\end{esercizio}

\begin{esercizio}[*]
\label{ese:6.48}
Risolvi i seguenti sistemi particolari.
\begin{enumeratea}
\item~\(\sistema{(y^2-4y+3)(x^2+2x-15)=0 
\\)x^2-3xy+2y^2)(9x^2-6xy+y^2)=0 }\)
\item~\(\sistema{(x-y)(x+4y-4)(x+y-1)(3x-5y-2)=0 
\\)3x+y-3)(x^2-4y^2)=0 }\)
\end{enumeratea}
\end{esercizio}

\paragraph{6.43.} 
a)~\((3;~-2)\vee(-3;~2)\vee\left(\dfrac{34} 7;~-\dfrac 1 7\right) 
\vee \left(-\dfrac{34} 7;~\dfrac 1 7\right)\),\quad 
b)~\(\emptyset\),\quad 
c)~\((-3;~0)\vee(3;~0)\),\protect\\ \quad 
d)~\(\left(\dfrac 1 2;~-2\right)\vee\left(-\dfrac 1 
2;~2\right) \vee \left(\dfrac 
7 4;~-\dfrac{11} 8\right)\vee\left(-\dfrac 7 4;~\dfrac{11} 8\right)\)

\paragraph{6.44.} 
a)~\((-2;~-1)\vee(2;~1)\),\quad 
b)~\((-1;~0)\vee(1;~0)\)

\paragraph{6.45.} 
a)~\((1;~1)\vee(3;~-3)\),\quad 
b)~\((3;~-1)\vee\left(\dfrac 
1 2;~\dfrac 
1 4\right)\),\quad c)~\((1;~0)\vee(-3;~-2)\),\protect\\
\quad d)~\((2;~0)\vee\left(2;~\dfrac 2 3\right)\vee(-6;~-2)\vee(-3;~1)\)

\paragraph{6.46.} a)~\((1;~-1)_\text{doppia} \vee(2;~0)\),\quad 
b)~\((0;~-1)_\text{doppia}\vee\left(-\dfrac 3 2;~-\dfrac 3 
2\right)\vee(1;~1)_\text{doppia}\),\quad c)~\((5;~8)\vee\left(\dfrac 3 
2;~1\right)\vee(-1;~-4)\vee\left(\dfrac 3 4;~-\dfrac 1 2\right)\),\quad 
d)~\((1;~-1)\vee(2;~0)\vee(-1;~1)\vee\left(\dfrac 1 2;~\dfrac 1 
2\right)\vee\left(-\dfrac 1 2;~-\dfrac 1 2\right)\vee\left(\dfrac{10} 
7;~-\dfrac 8 
7\right)\)

\paragraph{6.47.}a)~\((0;~0)_\text{doppia}\vee\left(3;~-\dfrac 3 
2\right)\vee(3;~\dfrac 3 4)\),\quad b)~\((t;~t)\vee\left(\dfrac 3 
2;~\dfrac 
1 
2\right)\vee\left(\dfrac 1 5;~-\dfrac 4 5\right)\vee\left(-\dfrac 1 
2;~-\dfrac 3 
2\right)\),\protect\\
\quad 
c)~\((0;~0)_\text{tripla}\vee(-2;~-2)_\text{doppia}\vee\left(-\dfrac 
2 
3,-\dfrac 4 3\right)_\text{doppia}\vee(6;~-2)\vee(-6;~-2)\),\quad 
d)~\(\left(\dfrac 3 
2;~\dfrac 3 2\right)\vee\left(-\dfrac 3 2;~\dfrac 3 2\right)\)

\paragraph{6.48.} 
a)~\((1;~1)\vee(2;~1)\vee(3;~3)_\text{doppia}\vee(6;~3)\vee\left(\dfrac 1 
3;~1\right)\vee(1;~3)\vee(-5;~-5)\vee\left(-5;~-\dfrac 5 
2\right)\vee\left(3;~\dfrac 3 
2\right)\vee(-5;~-15)\vee(3;~9)\),\quad 
b)~\((0;~0)_\text{doppia} \vee
\left(\dfrac 3 4;~\dfrac 3 4\right)\vee
(1;~0)\vee
\left(\dfrac {8}{11};~\dfrac {9}{11}\right)\vee
\left(\dfrac{17}{18};~\dfrac 1 6\right)\vee
\left(\dfrac 4 3;~\dfrac 2 3\right)\vee\protect\\
(-4;~2)(2;~-1)\vee
\left(\dfrac 2 3;~\dfrac 1 3\right)\vee
\left(\dfrac 4{11};~-\dfrac 2{11}\right)\vee
(4;~2)\)

\pagebreak %-------------------------------------
\subsection*{6.4 - Problemi che si risolvono con sistemi di grado 
superiore al 
primo}
\begin{htmulticols}{2}
\begin{esercizio}[*]
\label{ese:6.49}
La differenza tra due numeri è\(\dfrac {11} 4\) e il loro 
prodotto\(\dfrac 
{21} 8\) 
Trova i due numeri.
\end{esercizio}

\begin{esercizio}[*]
\label{ese:6.50}
Trovare due numeri positivi sapendo che la metà del primo supera di~\(1\) 
il 
secondo e che il quadrato del secondo supera di~\(1\) la sesta parte del 
quadrato del primo.
\end{esercizio}

\begin{esercizio}[*]
\label{ese:6.51}
Data una proporzione tra numeri naturali conosciamo i due medi che 
sono~\(5\) e 
\( 16\) Sappiamo anche che il rapporto tra il prodotto degli estremi e la 
loro 
somma è uguale a\(\dfrac {10} 3\) Trovare i due estremi.
\end{esercizio}

\begin{esercizio}[*]
\label{ese:6.52}
La differenza tra un numero di due cifre con quello che si ottiene 
scambiando le 
cifre è uguale a~\(36\) La differenza tra il prodotto delle cifre e la 
loro 
somma è uguale a~\(11\) Trovare il numero.
\end{esercizio}

\begin{esercizio}[*]
\label{ese:6.53}
Oggi la differenza delle età tra un padre e sua figlia è~\(26\) anni, 
mentre 
due 
anni fa il prodotto delle loro età era~\(56~\)Determina l'età del padre e 
della 
figlia.
\end{esercizio}

\begin{esercizio}[*]
\label{ese:6.54}
La somma delle età di due fratelli oggi è~\(46\) anni, mentre fra due 
anni 
la 
somma dei quadrati delle loro età sarà~\(1250\) Trova l'età dei due 
fratelli.
\end{esercizio}

\begin{esercizio}[*]
\label{ese:6.55}
Nella produzione di un oggetto la macchina A impiega 5 minuti in più 
rispetto 
alla macchina B. Determinare il numero di oggetti che produce ciascuna 
macchina 
in 8 ore se in questo periodo la macchina A ha prodotto 16 oggetti in 
meno 
rispetto alla macchina B.
\end{esercizio}

\begin{esercizio}[*]
\label{ese:6.56}
In un rettangolo la differenza tra i due lati è uguale a\(2\munit{cm}\) Se 
si 
diminuiscono entrambi i lati di~\(1\munit{cm}\) si ottiene un'area di 
\(0,1224\munit{m^2}~\)Calcolare il perimetro del rettangolo.
\end{esercizio}

\begin{esercizio}[*]
\label{ese:6.57}
Trova due numeri sapendo che la somma tra i loro quadrati è~\(100\) e il 
loro 
rapporto\( \dfrac 3 4\)
\end{esercizio}

\begin{esercizio}[*]
\label{ese:6.58}
Ho comprato due tipi di vino. In tutto 30 bottiglie. Per il primo tipo ho 
speso 
54 € e per il secondo 36 €. Il prezzo di una bottiglia del secondo tipo 
costa 
2,5 € in meno di una bottiglia del primo tipo. Trova il numero delle 
bottiglie 
di ciascun tipo che ho acquistato e il loro prezzo unitario.
\end{esercizio}

\begin{esercizio}[*]
\label{ese:6.59}
In un triangolo rettangolo di area\(630\munit{m^2}\), l'ipotenusa misura 
\(53\munit{m}~\)Determinare il perimetro.
\end{esercizio}

\begin{esercizio}[*]
\label{ese:6.60}
Un segmento di\(35\munit{cm}\) viene diviso in due parti. La somma dei 
quadrati 
costruiti su ciascuna delle due parti è\(625\munit{{cm}^2}\) Quanto misura 
ciascuna parte?
\end{esercizio}

\begin{esercizio}[*]
\label{ese:6.61}
Se in un rettangolo il perimetro misura~\(16,8\munit{m}\) e l'area\( 
17,28\munit{m^2}\), quanto misura la sua diagonale?
\end{esercizio}

\begin{esercizio}[*]
\label{ese:6.62}
In un triangolo rettangolo la somma dei cateti misura~\(10,5\munit{cm}\), 
mentre 
l'ipotenusa è~\(7,5\munit{cm}\) Trovare l'area.
\end{esercizio}

\begin{esercizio}[*]
\label{ese:6.63}
Quanto misura un segmento diviso in due parti, tali che una parte è\( 
\dfrac 3 4 
\) dell'altra, sapendo che la somma dei quadrati costruiti su ognuna 
delle 
due 
parti è uguale a\(121\munit{{cm}^2}\)?
\end{esercizio}

\begin{esercizio}[*]
\label{ese:6.64}
In un trapezio rettangolo con area di\(81\munit{m^2}\) la somma della base 
minore 
e dell'altezza è\(12\munit{m}\) mentre la base minore è\(\dfrac 1 5\) 
della 
base 
maggiore. Trovare il perimetro del rettangolo.
\end{esercizio}

\begin{esercizio}[*]
\label{ese:6.65}
La differenza tra le diagonali di un rombo è\(8\munit{cm}\), mentre la sua 
area è 
\(24\munit{{cm}^2}~\)Determinare il lato del rombo.
\end{esercizio}

\begin{esercizio}[*]
\label{ese:6.66}
Sappiamo che in un trapezio rettangolo con area di\(40\munit{{cm}^2}\) la 
base 
minore è\(7\munit{cm}\), mentre la somma della base maggiore e 
dell'altezza 
è 
\(17\munit{cm}\) Trovare il perimetro del rettangolo.
\end{esercizio}

\begin{esercizio}[*]
\label{ese:6.67}
Un rettangolo ha l'area uguale a quella di un quadrato. L'altezza del 
rettangolo 
è\(16\munit{cm}\), mentre la sua base è di\(5\munit{cm}\) maggiore del lato 
del 
quadrato. Determinare il lato del quadrato.
\end{esercizio}

\begin{esercizio}[*]
\label{ese:6.68}
La differenza tra i cateti di un triangolo rettangolo è\(7k\), mentre la 
sua 
area 
è\(60 k^2~\)Calcola il perimetro. (\(k>0\))
\end{esercizio}

\begin{esercizio}[*]
\label{ese:6.69}
L'area di un rettangolo che ha come lati le diagonali di due quadrati 
misura\(90 
k^2\) La somma dei lati dei due quadrati misura\(14k~\)Determinare i lati 
dei 
due 
quadrati. (\( k>0\))
\end{esercizio}

\begin{esercizio}[*]
\label{ese:6.70}
Nel rettangolo ABCD la differenza tra altezza e base è\(4k\) Se 
prolunghiamo 
la 
base AB dalla parte di B di\(2k\) fissiamo il punto E e congiungiamo B 
con 
E. 
Trovare il perimetro del trapezio AECD sapendo che la sua area è\(28k^2\) 
con 
\(k>0\)
\end{esercizio}

\begin{esercizio}[*]
\label{ese:6.71}
In un triangolo isoscele la base è\( \dfrac 2 3\) dell'altezza e l'area 
è\(12k^2\) 
Trova il perimetro del triangolo.
\end{esercizio}
\end{htmulticols}

\paragraph{6.49.}\(\left(-\dfrac 3 4;~-\dfrac 7 2\right)\vee \left(\dfrac 
7 
2;~-\dfrac 
3 4\right)\)
\paragraph{6.50.}\((12;~5)\)

\paragraph{6.51.}\((4;~20)\vee (20,4)\)

\paragraph{6.52.}\(73\)

\paragraph{6.53.}\((30;~4)\)

\paragraph{6.54.}\((23;~23)\)

\paragraph{6.55.}\((32;~48)\)

\paragraph{6.56.}\(2p=144\munit{cm}\)

\paragraph{6.57.}\((-6;~-8)\vee (6,8)\)

\paragraph{6.58.}\((12;~18)\)

\paragraph{6.59.}\(2p=126\munit{m}\)

\paragraph{6.60.}\([15\munit{cm}\) e\(20\munit{cm}]\)

\paragraph{6.61.}\(\text{Diagonale }=6\munit{m}\)

\paragraph{6.62.}\(\oarea=13,5\munit{{cm}^2}\)

\paragraph{6.63.}\(15,4\munit{cm}\)

\paragraph{6.64.}\(2p_1=42\vee 2p_2=57+3\sqrt{145}\)

\paragraph{6.65.}\(2\sqrt{10}\munit{cm}\)

\paragraph{6.66.}\(2p=24+2\sqrt{13}\)

\paragraph{6.67.}\(20\munit{cm}\)

\paragraph{6.68.}\(2p=40k\)

\paragraph{6.69.}\(l_1=5k\vee l_2=9k\)

\paragraph{6.70.}\(2p=15+k\sqrt{53}\)

\paragraph{6.71.}\(2p=4k(1+\sqrt{10})\)


